\par\Needspace{12\baselineskip}\bigskip
\pdfbookmark[2]{Exposition, Third section}{pallas-II-section3}
\begin{center}
{\large\bfseries Third section.\par}
\medskip
\textit{Method for calculating the finite variations in the elements over a limited period.}
\end{center}
\pdfbookmark[3]{Exposition, 17}{pallas-II-article17}
\begin{center}17.\end{center}
\noindent\hspace*{\parindent}The smallness of the perturbing masses compared with the mass of the Sun is an extremely favourable circumstance in the theory of planetary perturbations. Strictly speaking, in calculating the \SourceLeaf{467}instantaneous variations in each element, one would have to use the true values of the elements that enter into the expression for those variations. Given how small these variations themselves are, however, the result will differ very little if only approximate values of the elements are used. One obtains at least a first approximation in this way, which can be corrected later if deemed necessary.
\par First, then, one determines by the known methods purely elliptic elements that agree with the observations as closely as possible. Now, if only the theory of the planet's motion over a limited time is at issue, one begins by dividing this time into several equal intervals, each of which we shall take to be \(\theta\) days. It will be convenient to choose the first epoch of these intervals, which we shall denote by \(T\), a little before the period for which the motion is to be calculated, and to continue the series \(T,\ T+\theta,\ T+2\theta,\ T+3\theta\) etc. a little beyond that period. Using the formulas developed above and the approximate elements, one calculates the numerical values of the instantaneous variations in the different elements at the epochs \(T,\ T+\theta,\ T+2\theta\) etc., namely the values of \(\dfrac{\mathrm{d}i}{\mathrm{d}t},\ \dfrac{\mathrm{d}\AscendingNode}{\mathrm{d}t},\ \dfrac{\mathrm{d}n}{\mathrm{d}t},\ \dfrac{\mathrm{d}\varphi}{\mathrm{d}t},\ \dfrac{\mathrm{d}\tilde\omega}{\mathrm{d}t},\ \dfrac{\mathrm{d}\epsilon}{\mathrm{d}t}\). There is no need for \(\dfrac{\mathrm{d}a}{\mathrm{d}t}\), since the value of \(a\) is always linked to that of \(n\). It remains, then, to integrate these six differentials and to determine the mean motion \(\int n\mathrm{d}t\) by a second integration. The following method seems to us the most convenient for this purpose.
